%%%%%%%%%%%%%%%%%%%%%%%%%%%%%%%%%%%%%%%%%%%%%%%%%%%%%%%%%%
%   Autoren:
%   Prof. Dr. Bernhard Drabant
%   Prof. Dr. Dennis Pfisterer
%   Prof. Dr. Julian Reichwald
%%%%%%%%%%%%%%%%%%%%%%%%%%%%%%%%%%%%%%%%%%%%%%%%%%%%%%%%%%

%%%%%%%%%%%%%%%%%%%%%%%%%%%%%%%%%%%%%%%%%%%%%%%%%%%%%%%%%%
%	ANLEITUNG: 
%   1. Ersetzen Sie da Bild firmenlogo im Verzeichnis img
%   2. Passen Sie alle Stellen im Dokument an, die mit
%      @stud
%      markiert sind
%%%%%%%%%%%%%%%%%%%%%%%%%%%%%%%%%%%%%%%%%%%%%%%%%%%%%%%%%%

%%%%%%%%%%%%%%%%%%%%%%%%%%%%%%%%%%%%%%%%%%%%%%%%%%%%%%%%%%
%	ACHTUNG:
%   Für das Erstellen des Literaturverzeichnisses wird das
%   modernere Paket biblatex in Kombination mit biber
%   verwendet - nicht mehr das ältere Paket BibTex!
%
%   Bitte stellen Sie Ihre TeX-Umgebung entsprechend ein (z.B. TeXStudio):
%   Einstellungen --> Erzeugen --> Standard Bibliographieprogramm: biber
%%%%%%%%%%%%%%%%%%%%%%%%%%%%%%%%%%%%%%%%%%%%%%%%%%%%%%%%%%

\documentclass[fontsize=12pt,BCOR=5mm,DIV=12,parskip=half,listof=totoc,
    paper=a4,toc=bibliography,pointlessnumbers]{scrreprt}

%toc=listof,listof=entryprefix,

\makeindex

%% Elementare Pakete, Konfigurationen und Definitionen werden geladen (gegebenenfalls anpassen)
% !TEX root =  master.tex

%%%%%%%%%%%%%%%%%%%%%%%%%%%%%%%%%%%%%%%%%%%%%%%%%%%%%%%%%%%%%%%%%%
%	ANLEITUNG: 
% Passen Sie gegebenenfalls alle Stellen im Dokument an, die mit 
% @stud 
% markiert sind.
%%%%%%%%%%%%%%%%%%%%%%%%%%%%%%%%%%%%%%%%%%%%%%%%%%%%%%%%%%%%%%%%%%

%%
%% @stud
%%
%% LANGUAGE SETTINGS
\usepackage[ngerman]{babel} 	        % german language
\usepackage[german=quotes]{csquotes} 	% correct quoting using \enquote{}
%\usepackage[english]{babel}          % english language
%\usepackage{csquotes} 	              % correct quoting using \enquote{}

\usepackage{listings}	                %Format Listings properly
\usepackage{lipsum}                   % Blindtext
\usepackage{graphicx}                 % use various graphics formats
\usepackage[german]{varioref}         % nicer references \vref
\usepackage{caption}	                % better Captions
\usepackage{booktabs}                 % nicer Tabs
\usepackage[hidelinks=true]{hyperref} % keine roten Markierungen bei Links
\usepackage{fnpct}                    % Correct superscripts
\usepackage{calc}                     % Used for extra space below footsepline, in particular
\usepackage{array}
\usepackage{acronym}
\usepackage{algorithm}
\usepackage{algpseudocode}
\usepackage{setspace}
\usepackage{tocloft}

%% Schriftarten- und Zeichenpakete
\usepackage[T1]{fontenc}
\usepackage[utf8]{inputenc}

%%
%% @stud
%%
%%	FONT SELECTION: Schriftarten und Schriftfamilie
%%%%%%%%%%%%%
%% SCHRIFTART
%%%%%%%%%%%%%
% 0) without decomment: normal font families
% ...
% 1) Latin Modern
%\usepackage{lmodern}
% 2) Times
%\usepackage{mathptmx}
% 3) Helvetica
%\usepackage[scaled=.92]{helvet}
%%%%%%%%%%%%%%%%%%
%%	SCHRIFTFAMILIE
%%%%%%%%%%%%%%%%%%
% ohne Serifen
\renewcommand*{\familydefault}{\sfdefault}
\addtokomafont{disposition}{\sffamily}
%
% mit Serifen
%\renewcommand*{\familydefault}{\rmdefault}
%\addtokomafont{disposition}{\rmfamily}
%
% Typewriter
%\renewcommand*{\familydefault}{\ttdefault}
%\addtokomafont{disposition}{\ttfamily}

%%
%% @stud
%%
%% Uncomment the following lines to support hard URL breaks in bibliography
%\apptocmd{\UrlBreaks}{\do\f\do\m}{}{}
%\setcounter{biburllcpenalty}{9000}% Kleinbuchstaben
%\setcounter{biburlucpenalty}{9000}% Großbuchstaben

%%
%% @stud
%%
%% FOOTNOTES: Count footnotes over chapters
%% \counterwithout{footnote}{chapter}

%	ACRONYMS
\makeatletter
\@ifpackagelater{acronym}{2015/03/20}
{\renewcommand*{\aclabelfont}[1]{\textbf{{\acsfont{#1}}}}}{}
\makeatother

%	LISTINGS
% @stud: ggf. Namen/Text anpassen (englisch)
\renewcommand{\lstlistingname}{Quelltext}
\renewcommand{\lstlistlistingname}{Quelltextverzeichnis}
\lstdefinestyle{myStyle}{frame=tb}
\lstset{numbers=left,
	numberstyle=\tiny,
	captionpos=b,
	basicstyle=\ttfamily\small,
	style=myStyle}

%	ALGORITHMS
% @stud: ggf. Namen/Text anpassen (englisch)
\renewcommand{\listalgorithmname}{Algorithmenverzeichnis}
\floatname{algorithm}{Algorithmus}

%	PAGE HEADER / FOOTER
%	Warning: There are some redefinitions throughout the master.tex-file!  DON'T CHANGE THESE REDEFINITIONS!
\RequirePackage[automark]{scrlayer-scrpage}
%alternatively with separation lines: \RequirePackage[automark,headsepline,footsepline]{scrlayer-scrpage}

\renewcommand{\chaptermarkformat}{}
\RedeclareSectionCommand[beforeskip=0pt]{chapter}
\clearpairofpagestyles

%\ifoot[\rule{0pt}{\ht\strutbox+\dp\strutbox}DHBW Mannheim]{\rule{0pt}{\ht\strutbox+\dp\strutbox}DHBW Mannheim}
\ofoot[\rule{0pt}{\ht\strutbox+\dp\strutbox}\pagemark]{\rule{0pt}{\ht\strutbox+\dp\strutbox}\pagemark}
\ohead{\headmark}

\newcommand{\TitelDerArbeit}[1]{\def\DerTitelDerArbeit{#1}\hypersetup{pdftitle={#1}}}
\newcommand{\AutorDerArbeit}[1]{\def\DerAutorDerArbeit{#1}\hypersetup{pdfauthor={#1}}}
\newcommand{\Firma}[1]{\def\DerNameDerFirma{#1}}
\newcommand{\Kurs}[1]{\def\DieKursbezeichnung{#1}}
\newcommand{\Abteilung}[1]{\def\DerNameDerAbteilung{#1}}
\newcommand{\Studiengangsleiter}[1]{\def\DerStudiengangsleiter{#1}}
\newcommand{\WissBetreuer}[1]{\def\DerWissBetreuer{#1}}
\newcommand{\FirmenBetreuer}[1]{\def\DerFirmenBetreuer{#1}}
\newcommand{\Bearbeitungszeitraum}[1]{\def\DerBearbeitungszeitraum{#1}}
\newcommand{\Abgabedatum}[1]{\def\DasAbgabedatum{#1}}
\newcommand{\Matrikelnummer}[1]{\def\DieMatrikelnummer{#1}}
\newcommand{\Studienrichtung}[1]{\def\DieStudienrichtung{#1}}
\newcommand{\ArtDerArbeit}[1]{\def\DieArtDerArbeit{#1}}
\newcommand{\Literaturverzeichnis}{Literaturverzeichnis}

\newcommand{\settingBibFootnoteCite}{
	\setlength{\bibparsep}{\parskip}		  % Add some space between biblatex entries in the bibliography
	\addbibresource{bibliography.bib}	    % Add file bibliography.bib as biblatex resource
	\DefineBibliographyStrings{ngerman}{andothers = {{et\,al\adddot}},}
}

\newcommand{\setTitlepage}{
	% !TEX root =  master.tex
% @stud: ggf. Namen/Text anpassen (englisch)
\begin{titlepage}
	\begin{minipage}{\textwidth}
		\vspace{-2cm}
		\includegraphics[scale=0.6]{\imagedir/firmenlogo} \hfill \includegraphics{\imagedir/logo}
	\end{minipage}
	\vspace{1em}
%\sffamily
	\begin{center}
	{\textsf{\large Duale Hochschule Baden-Württemberg Mannheim}}\\[4em]
	{\textsf{\textbf{\large{\DieArtDerArbeit}arbeit}}}\\[6mm]
	{\textsf{\textbf{\Large{}\DerTitelDerArbeit}}} \\[1.5cm]
	{\textsf{\textbf{\large{}Studiengang Wirtschaftsinformatik}}\\[6mm]
	\textsf{\textbf{Studienrichtung \DieStudienrichtung}}}\vspace{10em}

	\begin{minipage}{\textwidth}
		\begin{tabbing}
			Wissenschaftliche(r) Betreuer(in): \hspace{0.85cm}\=\kill
			Verfasser(in): \> \DerAutorDerArbeit \\[1.5mm]
			Matrikelnummer: \> \DieMatrikelnummer \\[1.5mm]
			Firma: \> \DerNameDerFirma  \\[1.5mm]
			Kurs: \> \DieKursbezeichnung \\[1.5mm]
			Studiengangsleiter: \> \DerStudiengangsleiter \\[1.5mm]
			Bearbeitungszeitraum: \> \DerBearbeitungszeitraum\\[1.5mm]
		\end{tabbing}
	\end{minipage}
	\end{center}
\end{titlepage}
	\pagenumbering{roman} % Römische Seitennummerierung
	\normalfont
}

\newcommand{\initializeText}{
	\clearpage
	\ihead{\chaptername~\thechapter} % Neue Header-Definition
	\pagenumbering{arabic}           % Arabische Seitenzahlen
}

\newcommand{\initializeBibliography}{
	\ihead{}
	\printbibliography[title=\Literaturverzeichnis]
	\cleardoublepage
}

\newcommand{\initializeAppendix}{
	\appendix
	\ihead{}
	\cftaddtitleline{toc}{chapter}{Anhang}{}
}



%%
%% @stud
%%
%% PERSÖNLICHE ANGABEN (BITTE VOLLSTÄNDIG EINGEBEN zwischen den Klammern: {...})
%%
\ArtDerArbeit{Haus}
\TitelDerArbeit{Einzelne Handelsgeschäfte\nl Fracht- und Speditionsvertrag\nl Lager- und Transportgeschäfte}
\AutorDerArbeit{Juri Hofmann}
\Firma{ING DiBa AG}
\Kurs{WWI22SEB}
\Studienrichtung{Software Engineering}
\Matrikelnummer{4850943}
\Studiengangsleiter{Prof. Dr. Thomas Holey}
\Bearbeitungszeitraum{14.06.2023 -- 21.06.2023}

%%
%% @stud
%%
%% BIBLIOGRAPHY (@stud: Bibliographie-Stil wählen - Position und Indizierung)
%%  Auswahl zwischen:
%%   NUMERIC Style
%%   IEEE Style
%%   ALPHABETIC Style
%%   HARVARD Style
%%   CHICAGO Style
%%   (oder eigenen zulässigen Stil wählen)
%%
%%%%%%%%%%%%%
%% Zitierstil
%%%%%%%%%%%%%
% NUMERIC Style - e. g. [12]
\newcommand{\indextype}{numeric}
%
% IEEE Style - numeric kind of style
%\newcommand{\indextype}{ieee}
%
% ALPHABETIC Style - e. g. [AB12]
%\newcommand{\indextype}{alphabetic}
%
% HARVARD Style
%\newcommand{\indextype}{apa}
%
% CHICAGO Style
%\newcommand{\indextype}{authoryear}
%
%%%%%%%%%%%%%%%%%%%%%%
%% Position des Zitats
%%%%%%%%%%%%%%%%%%%%%%
\newcommand{\position}{inline}
%
% (!!) FOOTNOTE POSITION NOT RECOMMENDED IN MINT DOMAIN:
%\newcommand{\position}{footnote}

%% Final: Setzen des Zitierstils und der Zitatposition
\usepackage[backend=biber, autocite=\position, style=\indextype]{biblatex}
\settingBibFootnoteCite

%%
%% Definitionen und Commands
%%
\newcommand{\abs}{\par\vskip 0.2cm\goodbreak\noindent}
\newcommand{\nl}{\par\noindent}
\newcommand{\mcl}[1]{\mathcal{#1}}
\newcommand{\nowrite}[1]{}
\newcommand{\NN}{{\mathbb N}}
\newcommand{\imagedir}{img}

\makeindex

\begin{document}

    \setTitlepage
%%%%%%%%%%%%%%%%%%%%%%%%%%%%%%%%%%%
% VERZEICHNISSE und ABSTRACT
%
% @stud: ggf. nicht benötigte Verzeichnisse auskommentieren/löschen
%
    \tableofcontents
    \cleardoublepage

% Abkürzungsverzeichnis
% @stud: acronyms.tex bearbeiten
    % !TEX root =  master.tex
\clearpage
\chapter*{Abkürzungsverzeichnis}	
\addcontentsline{toc}{chapter}{Abkürzungsverzeichnis}

\begin{acronym}[XXXXXXX]
	\acro{ad}[AD]{Archiv für Diplomatik, Schriftgeschichte, Siegel- und Wappenkunde}
	\acro{BMBF}{Bundesministerium für Bildung und Forschung}	
	\acro{DHBW}{Duale Hochschule Baden-Württemberg}
	\acro{ecu}[ECU]{European Currency Unit}
	\acro{eu}[EU]{Europäische Union}
	\acro{RDBMS}[RDBMS]{Relational Database Management System}
\end{acronym}
    \cleardoublepage

%%%%%%%%%%%%%%%%%%%%%%%%%%%%%%%%%%%%%%%%%%%%%%%%%%%%%%%%%%%%%%%%%%%%%%%%%%%%%%%%%%%%%%%%%%
% KAPITEL UND ANHÄNGE
%
% @stud:
%   - nicht benötigte: auskommentieren/löschen
%   - neue: bei Bedarf hinzufügen mittels input-Kommando an entsprechender Stelle einfügen
%%%%%%%%%%%%%%%%%%%%%%%%%%%%%%%%%%%%%%%%%%%%%%%%%%%%%%%%%%%%%%%%%%%%%%%%%%%%%%%%%%%%%%%%%%

    \initializeText
    \onehalfspacing

%%%%%%%%%%%%%%%%%%%%%%%%%%%%%%%%%%%
% KAPITEL
%
% @stud: einzelne Kapitel bearbeiten und eigene Kapitel hier einfügen
%%%%%%%%%%%%%%%%%%%%%%%%%%%%%%%%%%%
    % !TEX root =  master.tex

\chapter[Einzelne Handelsgeschäfte]{Einzelne Handelsgeschäfte}

\section{Definition Handelskauf}
Der Handelskauf gemäß dem HGB ist eine besondere Form des Kaufvertrags, bei dem mindestens eine der beteiligten Parteien ein Kaufmann ist oder bei dem es sich um einen beiderseitigen Handelskauf handelt.

Ein Handelskauf liegt vor, wenn eine Vereinbarung über den Kauf von Waren getroffen wird. Dabei umfasst der Begriff "Waren" bewegliche Sachen, also Güter, die körperlich greifbar sind.


\section{Handelsgeschäft und beiderseitiges Handelsgeschäft}
Die Unterscheidung zwischen einem Handelsgeschäft und einem beiderseitigen Handelsgeschäft ist von Bedeutung, um die Anwendung spezifischer rechtlicher Vorschriften im Rahmen des Handelsrechts zu bestimmen. Gemäß § 1 Abs. 1 HGB wird ein Handelsgeschäft definiert als ein Geschäft, an dem mindestens eine der beteiligten Parteien ein Kaufmann ist. Ein Kaufmann ist eine natürliche oder juristische Person, die ein Handelsgewerbe betreibt.

Unter einem beiderseitigen Handelsgeschäft nach §§ 377-379 HGB hingegen handelt es sich um ein Geschäft, bei dem sowohl der Verkäufer als auch der Käufer Kaufleute sind. Das bedeutet, dass beide Parteien im Handelsregister eingetragen sind oder aufgrund anderer Kriterien als Kaufleute gelten. In solchen Fällen werden spezielle Vorschriften im HGB angewendet, die ausschließlich für beiderseitige Handelsgeschäfte gelten.

Die Unterscheidung zwischen einem Handelsgeschäft und einem beiderseitigen Handelsgeschäft ist wichtig, da sie die Rechtsstellung und die Pflichten der beteiligten Parteien beeinflusst. Im Falle eines beiderseitigen Handelsgeschäfts sind beide Parteien Kaufleute und unterliegen den besonderen Bestimmungen des HGB, die beispielsweise Regelungen zur Handelskaufmannschaft, Handelsbüchern und Handelsregistern enthalten. Für andere Handelsgeschäfte, bei denen nur eine Partei Kaufmann ist, gelten die allgemeinen Regelungen des HGB, die für sämtliche Handelsgeschäfte Anwendung finden.
\section{Anforderungen an einen Handelskauf}
Damit ein Handelskauf nach dem HGB vorliegt, ist es erforderlich, dass zumindest eine der beteiligten Parteien als Kaufmann im Sinne des HGB gilt. Die Kaufmannseigenschaft kann sich auf verschiedene Weisen ergeben und ist eng mit der gewerblichen Tätigkeit einer Person verbunden. Ein Kaufmann ist in der Regel eine natürliche oder juristische Person, die ein Handelsgewerbe betreibt.

Die Kaufmannseigenschaft kann aufgrund verschiedener Merkmale festgestellt werden. Eines dieser Merkmale ist beispielsweise die Eintragung im Handelsregister. Durch die Registrierung werden Kaufleute öffentlich bekannt gemacht und ihre Handelstätigkeiten werden transparent. Eine weitere Möglichkeit, die Kaufmannseigenschaft zu erlangen, besteht in der Mitgliedschaft in einer Handelskammer oder einer vergleichbaren Organisation, die den Status eines Kaufmanns bestätigt.

Die Kaufmannseigenschaft ist von großer Bedeutung, da sie Auswirkungen auf die rechtlichen Rahmenbedingungen des Handelskaufs hat. Das HGB enthält spezifische Vorschriften und Regelungen, die ausschließlich für Kaufleute gelten. Diese handelsrechtlichen Bestimmungen beinhalten beispielsweise Regelungen zur Buchführung, Bilanzierung, zum Handelsregister und zur Handelsfirma. Sie dienen der Schaffung von Rechtssicherheit und gewährleisten einheitliche Standards im Handelsverkehr.
\newpage\section{Geltungsbereich des Handelskaufs und Ausnahmen}
Der Geltungsbereich des Handelskaufs erstreckt sich ausschließlich auf den Erwerb von beweglichen Sachen, die als "Waren" bezeichnet werden. Diese Waren sind physische Objekte, die man anfassen und transportieren kann, wie beispielsweise Kleidung, Elektronikgeräte oder Maschinen. Im Handelskauf werden also spezifische Regelungen und Vorschriften für den Handel mit solchen Gegenständen festgelegt.

Es ist wichtig zu beachten, dass der Handelskauf keine Anwendung auf den Kauf von Grundstücken, Rechten oder Forderungen findet. Diese Vermögenswerte gelten nicht als Waren im Sinne des Handelskaufs. Der Handelskauf bezieht sich vielmehr auf den Austausch von materiellen Gütern. Dadurch wird sichergestellt, dass die spezifischen Regelungen des Handelsrechts nur für den Handel mit beweglichen Sachen gelten, während andere Rechtsbereiche für den Erwerb von Grundstücken oder immateriellen Vermögenswerten relevant sind. Somit sind die Bestimmungen des Handelskaufs nicht auf den Erwerb von Grundstücken oder anderen immateriellen Vermögenswerten anwendbar.

    % !TEX root =  master.tex

    \chapter[Speditionsvertrag]{Speditionsvertrag \footnote{Quelle:~\cite{Handelsrecht} \& \cite{ADSp} \& \cite{Speditionsbetriebslehre_und_Logistik}}}

    \section{Speditionsvertrag Definition}

    Ein Speditionsvertrag ist ein rechtlicher Vertrag zwischen einem Versender  und einem Spediteur, der die Beförderung von Gütern regelt. In diesem Vertrag verpflichtet sich der Spediteur, im Auftrag des Versenders den Transport von Waren oder Gütern von einem Ort zum anderen zu organisieren und durchzuführen.

    Der Speditionsvertrag ist ein spezieller Vertragstyp und unterliegt in vielen Ländern spezifischen rechtlichen Regelungen. Im deutschen Recht wird der Speditionsvertrag beispielsweise in den §§ 453 ff. des HGB geregelt.

    Im Speditionsvertrag werden verschiedene Aspekte der Transportleistung festgelegt, darunter:
    \begin{itemize}
        \item \textbf{Leistungsumfang:} Der Vertrag legt fest, welche Art von Transportleistung erbracht werden soll, z. B. Landverkehr, Seefracht oder Luftfracht. Außerdem werden die genauen Leistungen wie Abholung, Verpackung, Lagerung und Zustellung der Güter festgelegt.
        \item \textbf{Güterbeschreibung:} Der Vertrag enthält eine genaue Beschreibung der zu transportierenden Güter, einschließlich Art, Menge, Gewicht, Abmessungen, Verpackung usw.
        \item \textbf{Transportroute:} Der Vertrag kann die geplante Transportroute oder den vorgesehenen Transportweg angeben. Dies kann wichtige Informationen über den Ursprungsort, den Bestimmungsort und gegebenenfalls Zwischenstopps enthalten.
        \item \textbf{Haftung und Versicherung:} Der Speditionsvertrag regelt die Haftung des Spediteurs für Schäden oder Verluste während des Transports.
        \item \textbf{Vergütung:} Die Vergütung für die erbrachten Speditionsleistungen wird im Vertrag festgelegt.
        \item \textbf{Laufzeit und Kündigung:} Der Vertrag kann eine bestimmte Laufzeit haben oder für einen spezifischen Transportauftrag gelten. Es können auch Regelungen zur Kündigung oder Beendigung des Vertrags enthalten sein.

    Der Speditionsvertrag ist ein wichtiges Instrument, um die Rechte und Pflichten sowohl des Versenders als auch des Spediteurs zu regeln und einen reibungslosen Ablauf des Transportprozesses sicherzustellen.
\end{itemize}
    \section{Spediteur im Rahmen des Vertrages}
    Gemäß dem HGB wird der Spediteur im Rahmen des Speditionsvertrags als eine natürliche oder juristische Person definiert, die gewerbsmäßig Speditionsleistungen erbringt. Ein Spediteur kann ein eigenes Transportunternehmen betreiben oder als Vermittler zwischen dem Versender und den ausführenden Transportunternehmen fungieren.
    Der Spediteur übernimmt im Speditionsvertrag die Organisation und Durchführung des Transports im Namen des Versenders.
    Er handelt dabei im eigenen Namen, aber für Rechnung des Versenders.
    Der Spediteur wird somit als Dienstleister des Versenders beauftragt, den Transport der Güter von einem Ort zum anderen zu organisieren.
    Ein Spediteur kann ein eigenes Transportunternehmen betreiben oder als Vermittler zwischen dem Versender und den ausführenden Transportunternehmen fungieren.
    Der Spediteur übernimmt im Speditionsvertrag die Organisation und Durchführung des Transports im Namen des Versenders. Er handelt dabei im eigenen Namen, aber für Rechnung des Versenders.
    Der Spediteur wird somit als Dienstleister des Versenders beauftragt, den Transport der Güter von einem Ort zum anderen zu organisieren.

    Die Aufgaben und Verantwortlichkeiten des Spediteurs umfassen unter anderem:\nl
    \begin{itemize}
        \item \textbf{Auswahl der Transportmittel:}  Der Spediteur wählt die geeigneten Beförderungsmittel aus, um die Güter vom Abholort zum Bestimmungsort zu transportieren.Dies kann beispielsweise Lkw, Schiffe, Flugzeuge oder andere Transportmittel umfassen.
        \item \textbf{Verpackung und Kennzeichnung:} Der Spediteur sorgt dafür, dass die Güter ordnungsgemäß verpackt und gekennzeichnet sind, um einen sicheren Transport zu gewährleisten.
        \item \textbf{Transportorganisation:} Der Spediteur plant und organisiert den Ablauf des Transports, einschließlich der Koordination von Abholung, Lagerung, Umladung und Zustellung der Güter.
        \item \textbf{Frachtdokumente:} Der Spediteur erstellt und verwaltet die erforderlichen Frachtdokumente, wie beispielsweise Frachtbriefe oder Konnossemente.
        \item \textbf{Haftung:} Der Spediteur haftet für Schäden, Verluste oder Verzögerungen, die während des Transports auftreten, es sei denn, er kann nachweisen, dass diese auf Umstände zurückzuführen sind, die außerhalb seiner Kontrolle liegen.
    \end{itemize}

    Gemäß dem HGB ist der Spediteur ein eigenständiger Vertragspartner des Versenders und unterliegt spezifischen rechtlichen Regelungen im Zusammenhang mit dem Speditionsvertrag. 
    Es ist wichtig zu beachten, dass die genauen Bestimmungen und Haftungsregelungen im Speditionsvertrag selbst vereinbart werden können und somit von Fall zu Fall unterschiedlich sein können.
    
    \section{Unterschiede zwischen Speditions- und Frachtsvertrag}
    Es gibt einige Unterschiede zwischen dem Speditionsvertrag und dem Frachtvertrag, darunter fallen folgende:\nl

    \begin{itemize}
        \item \textbf{Vertragsparteien:} Im Speditionsvertrag sind die Vertragsparteien der Versender und der Spediteur, während im Frachtvertrag der Absender und der Frachtführer die Vertragspartner sind.
        \item \textbf{Leistungsumfang:} Im Speditionsvertrag organisiert der Spediteur den Transport im Namen des Versenders, während der Frachtführer im Frachtvertrag den tatsächlichen Transport der Güter durchführt.
        \item \textbf{Haftung:} Im Speditionsvertrag haftet der Spediteur für die ordnungsgemäße Organisation des Transports, während der Frachtführer im Frachtvertrag für den sicheren Transport der Güter verantwortlich ist. Die Haftungsbereiche können sich je nach Vertragstyp unterscheiden.
        \item \textbf{Durchführung des Transports:} Im Speditionsvertrag führt der Spediteur den Transport nicht selbst durch, sondern beauftragt möglicherweise andere Transportunternehmen oder Frachtführer. Im Frachtvertrag übernimmt der Frachtführer direkt den Transport der Güter.
        \item \textbf{Versicherung:} Im Speditionsvertrag kann der Spediteur Versicherungen für den Transport abschließen, um mögliche Risiken abzudecken. Im Frachtvertrag kann der Absender eine separate Transportversicherung abschließen, um die Güter während des Transports zu schützen.
        \item \textbf{Vertragszweck:} Der Speditionsvertrag hat den Zweck, die Organisation und Koordination des Transports sicherzustellen, während der Frachtvertrag darauf abzielt, den physischen Transport der Güter von einem Ort zum anderen zu gewährleisten.
    \end{itemize}

\section{Pflichten des Spediteurs im Rahmen des \ac{HGB}}

Gemäß dem Handelsgesetzbuch \ac{HGB} hat ein Spediteur im Rahmen des Speditionsvertrags verschiedene Pflichten und Verantwortlichkeiten. Hier sind die wichtigsten Pflichten, die ein Spediteur gemäß \ac{HGB} hat:

\begin{itemize}
    \item \textbf{Organisatorische Pflichten:} Der Spediteur ist verpflichtet, den Transport der Güter zu organisieren und die notwendigen Maßnahmen zu ergreifen, um einen reibungslosen Ablauf des Transports zu gewährleisten.
    \item \textbf{Sorgfaltspflicht:} Der Spediteur muss die erforderliche Sorgfalt walten lassen, um Schäden an den Gütern während des Transports zu vermeiden.
    \item \textbf{Informationspflicht:} Der Spediteur ist verpflichtet, den Versender über wichtige Informationen im Zusammenhang mit dem Transport zu informieren, wie beispielsweise über Verzögerungen, Schäden oder andere relevante Ereignisse.
    \item \textbf{Auswahl von Unterfrachtführern:} Der Spediteur kann unter Umständen Unterfrachtführer beauftragen, den Transport der Güter durchzuführen. Dabei ist der Spediteur verpflichtet, sorgfältig geeignete Unterfrachtführer auszuwählen und deren Zuverlässigkeit zu prüfen.
    \item \textbf{Einhaltung von Vorschriften:} Der Spediteur muss die gesetzlichen Vorschriften und behördlichen Bestimmungen im Zusammenhang mit dem Transport einhalten, wie beispielsweise Zollvorschriften, Sicherheitsvorschriften oder Transportgenehmigungen.
    \item \textbf{Rechnungslegung:} Der Spediteur ist verpflichtet, dem Versender eine ordnungsgemäße Rechnung über die erbrachten Speditionsleistungen auszustellen.
    \end{itemize}

Darüber hinaus kann der Speditionsvertrag weitere spezifische Pflichten enthalten, die zwischen den Vertragsparteien vereinbart wurden. Es ist wichtig, dass der Spediteur seine Pflichten gemäß dem Speditionsvertrag gewissenhaft erfüllt, um eine zuverlässige und professionelle Abwicklung des Transports sicherzustellen.

\section{Rechte des Spediteurs}
Im Rahmen des Speditionsvertrags gemäß dem HGB hat der Spediteur verschiedene Ansprüche gegenüber dem Versender. Diese Ansprüche dienen dazu, die Rechte und Interessen des Spediteurs zu schützen und eine angemessene Vergütung für die erbrachten Speditionsleistungen sicherzustellen. Hier sind die wichtigsten Ansprüche, die der Spediteur gemäß HGB gegenüber dem Versender geltend machen kann:
\begin{itemize}
    \item \textbf{Anspruch auf Vergütung:} Der Spediteur hat einen Anspruch auf die vereinbarte Vergütung für die erbrachten Speditionsleistungen. Die Höhe der Vergütung kann im Speditionsvertrag festgelegt werden oder auf Grundlage der gesetzlichen Bestimmungen erfolgen.
    \item \textbf{Anspruch auf Auslagenersatz:} Der Spediteur hat Anspruch auf Erstattung von angemessenen Auslagen, die im Zusammenhang mit der Durchführung des Transports entstanden sind. Dazu können beispielsweise Kosten für Transportmittel, Lagerung, Zölle oder Versicherungen gehören.
    \item \textbf{Anspruch auf Vorschuss:} In bestimmten Fällen, wie beispielsweise bei drohender Gefährdung des Spediteurs oder wenn der Versender zahlungsunfähig ist, hat der Spediteur das Recht, einen angemessenen Vorschuss auf die Vergütung zu verlangen.
    \item \textbf{Anspruch auf Zurückbehaltungsrecht:} Wenn der Versender seine Zahlungsverpflichtungen nicht erfüllt, kann der Spediteur unter bestimmten Bedingungen ein Zurückbehaltungsrecht ausüben und die Herausgabe der Güter verweigern, bis die offenen Zahlungen beglichen sind.
    \item \textbf{Anspruch auf Schadensersatz:} Falls der Versender seine vertraglichen Pflichten verletzt und dadurch dem Spediteur Schaden entsteht, kann der Spediteur einen Anspruch auf Schadensersatz geltend machen.
\end{itemize}

\section{Ergänzungen der derzeitigen Regeln}
Die Vorschriften des HGB bilden den grundlegenden Rahmen für die Bildung eines Speditionsvertrags zwischen zwei Parteien.

Sie legen die allgemeinen rechtlichen Prinzipien und Anforderungen fest, die für den Speditionsvertrag gelten.
Diese Vorschriften umfassen Regelungen zur Vertragsbildung, zur Leistungserbringung, zur Haftung, zur Vergütung und zu anderen relevanten Aspekten des Speditionsvertrags.

Jedoch können die Vorschriften des HGB allein nicht alle spezifischen Bedingungen und Details abdecken, die für einen konkreten Speditionsvertrag relevant sein können.
Um diese Lücke zu füllen und eine umfassendere Regelung zu ermöglichen, werden die Vorschriften des HGB durch bestimmte Bedingungen erweitert und ergänzt.

In Deutschland werden die ergänzenden Regelungen für den Speditionsvertrag durch die Allgemeinen Deutschen Spediteurbedingungen bereitgestellt.
Die ADSp sind ein standardisiertes Vertragswerk, das von verschiedenen Verbänden und Organisationen der Speditionsbranche entwickelt wurde.

Sie enthalten spezifische Bestimmungen und Konditionen, die auf die besonderen Anforderungen und Gegebenheiten des Speditionsvertrags zugeschnitten sind.
Die ADSp erweitern die Vorschriften des HGB, indem sie zusätzliche Regelungen zu verschiedenen Aspekten des Speditionsvertrags festlegen.

Diese können beispielsweise die Haftung des Spediteurs, die Versicherung der transportierten Güter, die Abwicklung von Zollformalitäten, die Zahlungsbedingungen oder andere vertragliche Bedingungen betreffen.
Durch die Einbeziehung der ADSp in den Speditionsvertrag können die Vertragsparteien auf ein etabliertes und anerkanntes Regelwerk zurückgreifen, das ihnen Klarheit und Rechtssicherheit bietet.

Es ist wichtig anzumerken, dass die ADSp nicht automatisch Teil des Speditionsvertrags sind.
Sie müssen ausdrücklich in den Vertrag aufgenommen werden, entweder durch ausdrückliche Vereinbarung oder durch Verweis auf die ADSp in den Vertragsunterlagen.

Die Parteien haben jedoch die Möglichkeit, die ADSp ganz oder teilweise zu übernehmen oder individuelle Vertragsbedingungen zu vereinbaren, solange diese nicht gegen zwingende gesetzliche Bestimmungen verstoßen.
Die Einbeziehung der ADSp in den Speditionsvertrag bietet den Vertragsparteien den Vorteil einer standardisierten und anerkannten Grundlage für ihre Rechte und Pflichten.
Sie erleichtert die Vertragsverhandlungen, minimiert das Risiko von Missverständnissen und Konflikten und erleichtert die effiziente Abwicklung des Speditionsvertrags.

\chapter[Frachtvertrag]{Frachtvertrag \footnote{Quelle:~\cite{Handelsrecht} \&~\cite{frachtvertrag}}}

\section[Definition]{Definition Frachtvertrag}

Ein Frachtvertrag ist ein rechtlicher Vertrag zwischen einem Absender und einem Frachtführer, der die Beförderung von Gütern regelt. Gemäß § 407 des \ac{HGB} wird der wesentliche Inhalt des Frachtvertrages definiert als die Hauptleistung, nämlich den Transport von Gütern. Der Absender übergibt dem Frachtführer die zu transportierenden Güter, und der Frachtführer verpflichtet sich, diese Güter an den vereinbarten Zielort zu befördern.

Ein Frachtvertrag kann sowohl für den Transport innerhalb einer Ortschaft als auch für den Transport über größere Distanzen abgeschlossen werden. Es gibt keine festgelegte Beschränkung bezüglich der verwendeten Transportmittel. Es können herkömmliche Transportmittel wie LKWs, Schiffe, Flugzeuge oder Züge verwendet werden, aber auch alternative Optionen wie Gepäckträger oder Dienstleute können in den Geltungsbereich eines Frachtvertrages fallen.

Im Frachtvertrag werden die Rechte und Pflichten der Vertragsparteien festgelegt, einschließlich der Haftung des Frachtführers für eventuelle Schäden oder Verluste der Güter während des Transports. Der Frachtvertrag unterliegt den Bestimmungen des \ac{HGB}, insbesondere den §§ 407 bis 452d, die spezifische Regelungen für Frachtverträge enthalten.

\section{Pflichten des Frachtführers}

Gemäß dem HGB sind dem Frachtführer im Rahmen des Frachtvertrags verschiedene Pflichten auferlegt. Diese umfassen:

\begin{itemize}
    \item \textbf{Betriebssichere Verladung (§ 412 HGB):} Der Frachtführer ist verantwortlich für die sichere und ordnungsgemäße Verladung der Güter, um Schäden während des Transports zu vermeiden.
    \item \textbf{Beförderung und Ablieferung innerhalb der Lieferfrist (§ 423 HGB):} Der Frachtführer muss die Güter innerhalb der vereinbarten oder üblichen Lieferfrist zum Empfänger befördern und dort rechtzeitig und unversehrt abliefern.
    \item \textbf{Befolgung von Weisungen (§ 418 und § 421 HGB):} Der Frachtführer muss nachträgliche Weisungen des Absenders und Empfängeranweisungen befolgen, sofern sie mit dem Frachtvertrag vereinbar sind.
    \item \textbf{Einholen von Weisungen bei Hindernissen (§ 419 HGB):} Falls während des Transports oder bei der Ablieferung Hindernisse auftreten, die den Vertrag beeinträchtigen könnten, muss der Frachtführer angemessene Weisungen einholen, um die Hindernisse zu überwinden.
    \item \textbf{Einziehen von Nachnahmen in bar (§ 422 HGB):} Wenn eine Nachnahme vereinbart wurde, hat der Frachtführer das Recht, den fälligen Betrag in bar bei der Ablieferung einzuziehen.
\end{itemize}

Diese Pflichten gewährleisten einen sicheren und reibungslosen Transport der Güter gemäß den Vereinbarungen des Frachtvertrags und dienen dem Schutz der Interessen sowohl des Absenders als auch des Empfängers.

\section{Rechte des Frachtführers}

Bei einem Frachtvertrag nach dem HGB hat der Frachtführer verschiedene Rechte, die ihm im Rahmen des Transports gewährt werden. Im Folgenden sind einige der wichtigsten Rechte aufgeführt:

\begin{itemize}
    \item \textbf{Leistungsverweigerungsrecht} Der Frachtführer hat gemäß § 418 HGB das Recht, die Güter zurückzuhalten, bis die vereinbarte Fracht oder andere fällige Zahlungen beglichen sind. Dies dient als Sicherheit für die Erfüllung der Zahlungsverpflichtungen des Versenders.
    \item \textbf{Zurückbehaltungsrecht:} Der Frachtführer hat gemäß § 418 HGB das Recht, die Ausführung des Frachtvertrags zu verweigern, wenn der Versender die vereinbarte Fracht oder andere fällige Zahlungen nicht geleistet hat. Der Frachtführer kann die Leistung bis zur vollständigen Begleichung der offenen Forderungen verweigern.L
\end{itemize}
Diese Rechte des Frachtführers nach dem Handelsgesetzbuch (HGB) dienen dazu, seine Interessen und Ansprüche im Rahmen des Frachtvertrags zu schützen. Sie geben dem Frachtführer die Möglichkeit, angemessene Entgelte für seine Dienstleistungen zu erhalten und eine Sicherheit für die Zahlungsforderungen zu haben. Gleichzeitig ermöglichen diese Rechte dem Frachtführer, seine Leistung einzubehalten oder zu verweigern, wenn die Zahlungsverpflichtungen nicht erfüllt werden.
\section{Frachtbrief}

Gemäß § 408 des \ac{HGB} ist der Frachtbrief eine spezielle Art von Dokument, das im Bereich des Frachtverkehrs verwendet wird. Der Frachtbrief erfüllt verschiedene Funktionen und dient als Beweisurkunde für den Abschluss und Inhalt des Frachtvertrages sowie für die Übernahme des Gutes in der angegebenen Beschaffenheit.

Der Frachtbrief enthält wichtige Informationen und Angaben über das transportierte Gut. Gemäß § 408 HGB wird der Frachtbrief als Nachweis für folgende Punkte verwendet, bis zum Beweis des Gegenteils:

\begin{itemize}
    \item Den äußerlich guten Zustand des Gutes bei Übernahme durch den Frachtführer.
    \item Die Richtigkeit der Angaben über Anzahl der Frachtstücke, Zeichen und Nummern.
    \item Das überprüfte Gewicht, die Menge oder der Inhalt, sofern vom Frachtführer bescheinigt.
\end{itemize}

Der Frachtbrief hat eine hohe rechtliche Bedeutung, da er als Beweismittel verwendet werden kann. Er ermöglicht es den beteiligten Parteien, im Streitfall ihre Rechte und Pflichten im Zusammenhang mit dem Frachtvertrag nachzuweisen. Der Frachtbrief trägt zur Transparenz und Sicherheit im Frachtverkehr bei und erleichtert die Abwicklung von Frachtgeschäften.

\section{Haftung der beiden Parteien}
Gemäß dem HGB gibt es Haftungsgrundlagen sowohl für den Frachtführer als auch für den Absender im Rahmen eines Frachtvertrages.

Für den Frachtführer gibt es folgende Haftungsgrundlagen:
\begin{itemize}
    \item \textbf{Obhutshaftung:} Der Frachtführer haftet für Verlust, Beschädigung oder Verspätung des Transportguts, da er während des Transports die Obhut über das Gut hat.
    \item \textbf{Verschuldenshaftung:} Der Frachtführer haftet, wenn er eigenes Verschulden nachweisen kann, das zu Schäden am Transportgut geführt hat. Dies umfasst Fahrlässigkeit oder vorsätzliches Handeln.
    \item \textbf{Haftungsbefreiung und Haftungsbegrenzung:} Es können bestimmte Umstände eintreten, die den Frachtführer von seiner Haftung befreien oder seine Haftung begrenzen. Zum Beispiel können Naturereignisse, Handlungen des Absenders oder unvorhersehbare Umstände die Haftung des Frachtführers beschränken.
    \item \textbf{Haftung für Handlungen von Hilfspersonen:} Der Frachtführer kann auch für Handlungen oder Verschulden von Personen haftbar gemacht werden, die er zur Erfüllung des Frachtvertrags hinzuzieht, wie beispielsweise Unterauftragnehmer oder Mitarbeiter.
\end{itemize}

Der Absender haftet in der Regel nicht für Schäden am Transportgut, da seine Verpflichtungen hauptsächlich darauf abzielen, das Gut ordnungsgemäß zu verpacken, zu kennzeichnen und die erforderlichen Begleitpapiere bereitzustellen. Jedoch haftet der Absender dem Frachtführer für die Zahlung der vereinbarten Fracht und der erforderlichen Aufwendungen gemäß § 420 \ac{HGB}.
    % !TEX root =  master.tex


\chapter[Lagergeschäfte]{Lagergeschäfte}

\section{Definition}
Lagergeschäfte im Sinne des HGB beziehen sich auf eine spezielle Art von entgeltlichen Verträgen, die zwischen einem Lagerhalter und einem Einlagerer abgeschlossen werden. Diese Verträge regeln die Verwahrung und Aufbewahrung von Waren durch den Lagerhalter.

Gemäß § 467 HGB liegt ein Lagergeschäft vor, wenn der Lagerhalter die Verantwortung für die ordnungsgemäße Aufbewahrung, Pflege und Sicherheit der eingelagerten Güter übernimmt. Dabei handelt es sich um eine Dienstleistung, bei der der Lagerhalter gegen Entgelt die Waren des Einlagerers in seinem Lager aufnimmt und für deren Verwahrung sorgt.

Ein Lagergeschäft kann verschiedene Formen annehmen. Es kann sich beispielsweise um die Lagerung von Rohstoffen, Halbfertigprodukten oder Endprodukten handeln. Die Waren können sowohl unbewegliche als auch bewegliche Güter umfassen. Typische Beispiele für Lagergeschäfte sind die Lagerung von Waren in Lagerräumen, Lagerhallen oder anderen Lagerstätten.

Die Definition von Lagergeschäften im HGB legt den Fokus auf die rechtliche und vertragliche Beziehung zwischen dem Lagerhalter und dem Einlagerer. Der Lagerhalter übernimmt die Verantwortung für die ordnungsgemäße Verwahrung der Waren und hat die Pflicht, diese vor Schäden oder Verlusten zu schützen. Der Einlagerer hingegen vertraut dem Lagerhalter die sichere Aufbewahrung seiner Waren an.
\section{Rechte und Pflichten des Lagerhalters}

Die Rechte und Pflichten des Lagerhalters sind im HGB geregelt und beziehen sich auf seine Verantwortung für die ordnungsgemäße Verwahrung und Aufbewahrung der eingelagerten Waren. Als Lagerhalter übernimmt er eine wichtige Rolle und hat verschiedene Aufgaben und Verpflichtungen zu erfüllen.
\begin{itemize}
    \item \textbf{Aufbewahrungspflicht:} Der Lagerhalter ist verpflichtet, die ihm anvertrauten Waren sorgfältig und ordnungsgemäß aufzubewahren. Er muss dafür sorgen, dass die Waren vor Schäden, Verlusten oder Diebstahl geschützt sind.
    \item \textbf{Bestandsaufnahme und Dokumentation:} Der Lagerhalter hat die Pflicht, eine genaue Bestandsaufnahme der eingelagerten Waren vorzunehmen und diese zu dokumentieren. Dies beinhaltet die Erfassung von Art, Menge, Zustand und Wert der Waren.
    \item \textbf{Herausgabepflicht:} Der Lagerhalter ist verpflichtet, die eingelagerten Waren gemäß den Anweisungen des Einlagerers herauszugeben. Der Lagerhalter muss die Waren ordnungsgemäß verpackt und unversehrt übergeben.
    \item \textbf{Vergütung:} Der Lagerhalter hat Anspruch auf eine angemessene Vergütung für seine Lagerdienstleistung.
    \item \textbf{Pfandrecht:} Der Lagerhalter besitzt ein gesetzliches Pfandrecht an den eingelagerten Waren. Dies bedeutet, dass er unter bestimmten Voraussetzungen ein Zurückbehaltungsrecht ausüben kann, um offene Forderungen gegenüber dem Einlagerer zu sichern. Das Pfandrecht ermöglicht es dem Lagerhalter, die Herausgabe der Waren zu verweigern, bis die offenen Zahlungen beglichen sind.
\end{itemize}

\section{Haftung des Lagerhalters für Schäden oder Verluste}
Die Haftung des Lagerhalters ist ein wesentlicher Aspekt im Zusammenhang mit Lagergeschäften. Sie bezieht sich auf die Verantwortung des Lagerhalters für Schäden oder Verluste, die während der Lagerung der Waren auftreten können. Im HGB sind spezifische Bestimmungen zur Haftung des Lagerhalters festgelegt, die sowohl seine Rechte als auch seine Pflichten regeln.

\begin{itemize}
    \item \textbf{Haftungsumfang:} Der Lagerhalter haftet für Schäden oder Verluste der Waren, außer er kann nachweisen, dass er seine Sorgfaltspflicht erfüllt hat und der Schaden unabhängig von seinem Verschulden eingetreten ist.
    \item \textbf{Sorgfaltspflicht:} Der Lagerhalter muss die Waren sorgfältig behandeln, indem er angemessene Lagerbedingungen sicherstellt und regelmäßige Kontrollen durchführt.
    \item \textbf{Beschränkung der Haftung:} Der Lagerhalter kann seine Haftung im Lagervertrag beschränken oder ausschließen, sofern dies klar und eindeutig vereinbart wird und keine gesetzlichen Bestimmungen verletzt werden.
    \item \textbf{Versicherung:} Der Lagerhalter kann eine Versicherung abschließen, um die Haftung zu begrenzen und mögliche Risiken abzudecken.
    \item \textbf{Ausnahme von der Haftung:} Der Lagerhalter kann von seiner Haftung befreit sein, wenn der Schaden oder Verlust auf Umständen höherer Gewalt beruht.
    \item \textbf{Schadensersatzansprüche:} Einlagerer haben das Recht auf Schadensersatz, der den Wert der beschädigten oder verlorenen Waren sowie mögliche Folgeschäden umfasst.
\end{itemize}

Die Rechte und Pflichten des Lagerhalters sind im HGB geregelt und beziehen sich auf seine Verantwortung für die ordnungsgemäße Verwahrung und Aufbewahrung der eingelagerten Waren. Als Lagerhalter übernimmt er eine wichtige Rolle und hat verschiedene Aufgaben und Verpflichtungen zu erfüllen.

\section{Kündigung und Beendigung des Lagervertrags}
Die Kündigung und Beendigung des Lagervertrags ist ein wichtiger Aspekt im Bereich der Lagergeschäfte. Eine korrekte Kündigung und Beendigung des Lagervertrags gemäß den Bedingungen und gesetzlichen Vorschriften ist wichtig, um einen reibungslosen Abschluss der Lagervereinbarung zu gewährleisten.

\begin{itemize}
    \item \textbf{Kündigung durch den Lagerhalter:} Der Lagerhalter kann den Lagervertrag gemäß den vereinbarten Bedingungen und Fristen schriftlich kündigen.
    \item \textbf{Kündigung durch den Einlagerer:} Der Einlagerer kann den Lagervertrag ebenfalls schriftlich kündigen, unter Beachtung der vereinbarten Kündigungsbedingungen.
    \item \textbf{Beendigung durch Erfüllung des Vertrags:} Der Lagervertrag endet automatisch, wenn die vereinbarte Lagerdauer abgelaufen ist und die Waren ordnungsgemäß zurückgegeben wurden.
    \item \textbf{Außerordentliche Kündigung:} In Ausnahmefällen kann eine außerordentliche Kündigung gerechtfertigt sein, wenn eine Vertragspartei wesentliche Pflichten verletzt oder eine Fortsetzung unzumutbar wäre.
    \item \textbf{Abwicklung der Beendigung:} Bei Vertragsbeendigung müssen sowohl der Lagerhalter als auch der Einlagerer ihre Verpflichtungen erfüllen, wie die ordnungsgemäße Warenübergabe und Zahlung offener Beträge.
\end{itemize}





\chapter[Transportgeschäfte]{Transportgeschäfte}

\section{Definition}

Transportgeschäfte sind Verträge, bei denen der Transport von Gütern gegen Entgelt erfolgt. Sie umfassen den Land-, Wasser- und Lufttransport, können national oder international sein und werden von Frachtführern, Speditionen oder Logistikunternehmen durchgeführt. Der Kern eines Transportgeschäfts besteht aus der Vereinbarung zwischen Absender und Frachtführer über den Ort, die Zeit und die Bedingungen des Transports. Es fallen Frachtkosten an, und der Frachtführer haftet gemäß dem HGB für Verluste oder Beschädigungen der Güter. Die sorgfältige Planung und Durchführung von Transportgeschäften gewährleistet einen reibungslosen Warentransport. Das HGB bietet den rechtlichen Rahmen für den Transport von Gütern und legt die Rechte und Pflichten der Beteiligten fest.


\section{Beförderungsvertrag}
Ein Beförderungsvertrag ist ein rechtlicher Vertrag, der zwischen dem Beförderer und dem Versender abgeschlossen wird und den Transport von Gütern regelt. Es handelt sich um einen zentralen Bestandteil des Transportwesens, da er die rechtlichen Grundlagen für den Transport von Waren schafft und die Rechte und Pflichten der beteiligten Parteien festlegt.

Im Beförderungsvertrag werden verschiedene Elemente definiert. Zunächst werden die Parteien des Vertrags festgelegt. Der Beförderer ist die Person oder das Unternehmen, das den Transport der Güter durchführt, während der Versender die Partei ist, die die Güter zur Beförderung übergibt.

Der Beförderungsvertrag legt auch die Art der zu befördernden Güter fest. Dabei kann es sich um physische Güter wie Waren, Rohstoffe oder Maschinen handeln, aber auch um immaterielle Güter wie Daten oder Informationen, die elektronisch übertragen werden.

Ein weiterer wichtiger Bestandteil des Beförderungsvertrags ist die Festlegung des Transportwegs und des Transportmittels. Der Vertrag kann beispielsweise die Route von einem bestimmten Startpunkt zum Zielpunkt vorgeben und das gewählte Transportmittel spezifizieren, wie beispielsweise Straßentransport, Schifffahrt, Luftfracht oder Schienentransport.

Im Beförderungsvertrag werden auch die Bedingungen und Modalitäten des Transports geregelt.
\section{Haftung des Frachtführers für Schäden oder Verluste}
Der Frachtführer ist die Person oder das Unternehmen, das den Transport von Gütern im Rahmen eines Frachtvertrags durchführt. Gemäß dem HGB und anderen einschlägigen Gesetzen haftet der Frachtführer für Schäden, die während des Transports an den beförderten Gütern entstehen.

Die Haftung des Frachtführers basiert auf dem Prinzip der Verschuldenshaftung. Das bedeutet, dass der Frachtführer für Schäden oder Verluste haftet, die auf sein eigenes Verschulden, das Verschulden seiner Mitarbeiter oder auf Mängel in seiner Organisation oder Ausrüstung zurückzuführen sind. Die Haftung des Frachtführers besteht unabhängig davon, ob der Schaden durch Fahrlässigkeit, Vorsatz oder andere Umstände verursacht wurde.

Die Haftung des Frachtführers kann in verschiedene Phasen des Transports unterteilt werden:

\begin{itemize}
    \item \textbf{Vorannahme:} Der Frachtführer übernimmt die Güter und prüft ihre Beförderungstauglichkeit.
    \item \textbf{Beförderung:} Der Frachtführer ist für die sichere und termingerechte Lieferung der Güter verantwortlich.
    \item \textbf{Ablieferung:} Der Frachtführer übergibt die Güter an den Empfänger oder den berechtigten Empfangsberechtigten.
\end{itemize}

Die Haftung des Frachtführers erstreckt sich über den gesamten Transportprozess und umfasst die Verantwortung für Schäden oder Verluste, die auf sein Verschulden, das Verschulden seiner Mitarbeiter oder auf Mängel in seiner Ausrüstung oder Organisation zurückzuführen sind.


    % !TEX root =  master.tex
\chapter{Zusammenfassung}

Im Laufe der Arbeit wurde deutlich, dass das Handelsrecht im HGB eine bedeutende Rolle für die Regelung und den Schutz von Geschäften im Handelsverkehr spielt.

Persönlich finde ich diese Themen äußerst relevant und praxisnah. Das Verständnis der einzelnen Handelsgeschäfte nach dem HGB ermöglicht es den Akteuren im Handelsverkehr, ihre Rechte und Pflichten besser zu verstehen und angemessen zu handeln. Die Regelungen zum Speditions- und Frachtvertrag bieten eine klare rechtliche Grundlage für den Transport von Gütern und stellen sicher, dass sowohl der Spediteur als auch der Frachtführer bestimmten Haftungsbestimmungen unterliegen.

Die intensive Auseinandersetzung mit diesen Themen hat mir ein besseres Verständnis für die rechtlichen Aspekte des Handelsverkehrs vermittelt. Die Kenntnis dieser rechtlichen Bestimmungen ist von entscheidender Bedeutung, um potenzielle Risiken zu minimieren und rechtliche Konflikte zu vermeiden.

Die Hausarbeit hat mir gezeigt, dass das HGB ein unverzichtbares Instrument ist, um den Handelsverkehr zu regeln und einen reibungslosen Ablauf von Handelsgeschäften zu gewährleisten. Außerdem war es ebenfalls eine große Hilfe bei der korrekten Recherche in dieser Hausarbeit. Es ist von großer Bedeutung, dass Unternehmen und Einzelpersonen, die im Handelsverkehr tätig sind, sich mit den Bestimmungen des HGB vertraut machen und ihre Handelsaktivitäten entsprechend gestalten.

Insgesamt bin ich der Überzeugung, dass die Themen der einzelnen Handelsgeschäfte, des Speditions- und Frachtvertrags sowie der Lager- und Transportgeschäfte nach dem HGB eine solide Grundlage für den Handelsverkehr bilden. Sie bieten Rechtssicherheit, ermöglichen effiziente Geschäftsabläufe und schützen die Interessen der Vertragsparteien. Daher sollten Unternehmen und Einzelpersonen, die im Handelsverkehr tätig sind, sich intensiv mit diesen Themen auseinandersetzen, um von den rechtlichen Vorteilen und Schutzmechanismen des HGB zu profitieren.




%%%%%%%%%%%%%%%%%%%%%%%%%%%%%%%%%%%
% ANHÄNGE
%
% @stud: einzelne Anhänge bearbeiten und eigene Anhänge hier einfügen
%        die nachfolgenden Zeilen deaktivieren, wenn keine Anhänge verwendet werden
%
%\initializeAppendix

%%%%%%%%%%%%%%%%%%%%%%%%%%%%%%%%%%%

    \singlespacing

%%%%%%%%%%%%%%%%%%%%%%%%%%%%%%%%%%%
% LITERATURVERZEICHNIS
% @stud: Literaturverzeichnis in Datei bibliography.bib anpassen.
%
% Alternative zu Verwendung von \initializeBibliography: Citavi ...
% (dann \initializeBibliography auskommentieren und eigenes LaTex Coding verwenden)
%
    \initializeBibliography
%%%%%%%%%%%%%%%%%%%%%%%%%%%%%%%%%%%

\end{document}
